\section{Quotients of seminorm spectra}\label{sec:general-spectrum-quotients}

Let $A$ be a commutative ring and let $\mathcal S(A)$ denote its space of
real-valued multiplicative ultrametric seminorms, with the evaluation topology.
For the relative statements, let $R$ be a commutative ring equipped with a fixed
seminorm $v$, let $A$ be a commutative $R$-algebra, and let $\mathcal S_v(A)$
be the subspace of seminorms restricting to $v$ on $R$.
No field, trivial valuation, finite generation, or reducedness assumption is imposed.

\begin{definition}[Zero locus in a seminorm spectrum]\label{spectrum:quotient-zero-locus}
\lean{EnvelopeProperty.Berkovich.Spectrum.zeroLocus,
EnvelopeProperty.Berkovich.RelativeSpectrum.zeroLocus}
\leanok
\uses{spectrum:affine-seminorm,spectrum:relative}
For an ideal $I\subseteq A$, set
$V_{\mathcal S}(I)=\{x\in\mathcal S(A):I\subseteq\operatorname{supp}(x)\}$.
In a relative spectrum, use the same condition to define $V_{\mathcal S_v}(I)$.
Equivalently, the seminorm must vanish on every element of $I$.
\end{definition}

\begin{lemma}[Closed zero loci]\label{spectrum:quotient-zero-locus-closed}
\lean{EnvelopeProperty.Berkovich.Spectrum.isClosed_zeroLocus,
EnvelopeProperty.Berkovich.RelativeSpectrum.isClosed_zeroLocus}
\leanok
\uses{spectrum:quotient-zero-locus,spectrum:topology}
Both $V_{\mathcal S}(I)$ and $V_{\mathcal S_v}(I)$ are closed in their respective spectra.
\end{lemma}
\begin{proof}
For each $a\in I$, the condition $|a|_x=0$ is closed because evaluation at $a$
is continuous and $\{0\}$ is closed in $\R_{\ge0}$. The asserted zero locus is
the intersection of these conditions. In the relative spectrum the same argument
uses the restriction of each evaluation map to the prescribed-base subspace.
\end{proof}

\begin{theorem}[Quotient spectrum]\label{spectrum:quotient-homeomorphism}
\lean{EnvelopeProperty.Berkovich.Spectrum.quotientHomeomorph,
EnvelopeProperty.Berkovich.Spectrum.isClosedEmbedding_quotientComap,
EnvelopeProperty.Berkovich.RelativeSpectrum.quotientHomeomorph,
EnvelopeProperty.Berkovich.RelativeSpectrum.isClosedEmbedding_quotientComap}
\leanok
\uses{spectrum:quotient-zero-locus-closed,spectrum:pullback,spectrum:relative}
Restriction along $A\to A/I$ identifies $\mathcal S(A/I)$ homeomorphically
with $V_{\mathcal S}(I)$ and identifies $\mathcal S_v(A/I)$ homeomorphically
with $V_{\mathcal S_v}(I)$. The restriction maps are closed embeddings.
For the inverse, $|a+I|=|a|_x$.
\end{theorem}
\begin{proof}
A seminorm pulled back from $A/I$ vanishes on $I$. Conversely, suppose $x$ vanishes
on $I$. If $a-b\in I$, then the ultrametric inequality gives
$|a|_x\le\max\{|b|_x,|a-b|_x\}=|b|_x$ and, by interchanging $a$ and $b$,
$|b|_x\le|a|_x$. Thus $|a+I|=|a|_x$ is well defined.
The multiplicative and ultrametric identities descend along the quotient map,
as do the values at zero and one. The two constructions are inverse because
every quotient class has a representative. They preserve a prescribed seminorm
on $R$, since the image of $r\in R$ in $A/I$ is the quotient class of its image
in $A$. Restriction is continuous by the evaluation definition of the topology.
Every evaluation on the inverse, at a quotient class represented by $a$, is
the restriction of evaluation at $a$ on $\mathcal S(A)$, and hence is continuous.
The same proof applies on the relative subspaces. The images are closed by
the preceding lemma, so both restriction maps are closed embeddings.
\end{proof}

\begin{theorem}[Surjective morphisms of rings and algebras]\label{spectrum:surjective-closed-embedding}
\lean{EnvelopeProperty.Berkovich.Spectrum.isClosedEmbedding_comap_of_surjective,
EnvelopeProperty.Berkovich.RelativeSpectrum.isClosedEmbedding_comap_of_surjective}
\leanok
\uses{spectrum:quotient-homeomorphism,spectrum:pullback}
A surjective ring homomorphism $A\to B$ induces a closed embedding
$\mathcal S(B)\to\mathcal S(A)$. A surjective $R$-algebra homomorphism
induces a closed embedding $\mathcal S_v(B)\to\mathcal S_v(A)$.
\end{theorem}
\begin{proof}
The first isomorphism theorem identifies $B$ with the quotient of $A$ by the
kernel of the homomorphism; for an $R$-algebra homomorphism this identification
is an $R$-algebra isomorphism. An isomorphism induces mutually inverse
continuous pullback maps on seminorm spectra, since every evaluation pulls back
to evaluation on the image of the element. The same is true on the relative
subspaces for an $R$-algebra isomorphism. The desired maps consequently factor
as one of these homeomorphisms followed by the closed quotient embedding.
\end{proof}

\begin{theorem}[Invariance under reduction]\label{spectrum:reduction-homeomorphism}
\lean{EnvelopeProperty.Berkovich.Spectrum.reductionHomeomorph,
EnvelopeProperty.Berkovich.RelativeSpectrum.reductionHomeomorph}
\leanok
\uses{spectrum:quotient-homeomorphism}
If $A_{\mathrm{red}}=A/\sqrt{(0)}$, restriction gives homeomorphisms
$\mathcal S(A_{\mathrm{red}})\simeq\mathcal S(A)$ and
$\mathcal S_v(A_{\mathrm{red}})\simeq\mathcal S_v(A)$.
They preserve the value of every regular function and its quotient class.
\end{theorem}
\begin{proof}
The support of any multiplicative seminorm is prime: it is an ideal by the
ultrametric inequality and multiplicativity, is proper since $|1|=1$, and
$|ab|=0$ implies $|a|=0$ or $|b|=0$. Every prime ideal contains the nilradical,
so every seminorm vanishes on $\sqrt{(0)}$. The zero loci in the quotient
homeomorphism theorem are therefore the entire respective spectra, and the
inverse evaluation formula proves compatibility with all regular functions.
\end{proof}

\begin{remark}\label{spectrum:quotient-specialization}
\uses{spectrum:quotient-homeomorphism,spectrum:surjective-closed-embedding,spectrum:reduction-homeomorphism}
The affine closed-immersion and reduction statements for a trivially valued field
are the special cases in which $R$ is that field and the prescribed base seminorm
is trivial. These general results concern the evaluation topology on real-valued
seminorms of arbitrary rings. They impose no boundedness or completeness condition
and do not assert the construction of the analytic spectrum of a Banach ring.
\end{remark}
